\documentclass[11pt,a4paper]{article}
\usepackage[utf8]{inputenc}
\usepackage{amsmath,amssymb}
\usepackage{geometry}
\geometry{margin=2.5cm}
\usepackage{listings}
\usepackage{xcolor}
\usepackage{hyperref}

\lstset{
  language=C++,
  basicstyle=\ttfamily\small,
  keywordstyle=\color{blue}\bfseries,
  commentstyle=\color{green!60!black},
  stringstyle=\color{red!70!black},
  numbers=left,
  numberstyle=\tiny\color{gray},
  numbersep=8pt,
  frame=single,
  breaklines=true,
  tabsize=4,
  showstringspaces=false
}

\title{IOI 1997 -- The Buses}
\author{Solution}
\date{}

\begin{document}
\maketitle

\section{Problem Statement}

A person observes bus arrivals at a stop and records arrival times during an observation
period. From these times, determine the bus routes (schedules).

Each bus route is characterized by:
\begin{itemize}
  \item A \textbf{start time} $s$ (first arrival).
  \item An \textbf{interval} $d > 0$ (time between consecutive buses).
  \item A \textbf{count} $c \ge 2$ (number of buses on this route).
\end{itemize}

Every recorded arrival must belong to exactly one route. Find a set of routes that
explains all arrivals using the \textbf{minimum number of routes}.

\medskip
\noindent\textbf{Constraints:} Number of arrivals $n \le 300$, arrival times $\le 59$.

\section{Solution Approach}

\subsection{Backtracking with Pruning}

The small time range (0--59) makes backtracking feasible:

\begin{enumerate}
  \item Maintain a frequency array $\mathtt{cnt}[0..59]$ of unassigned arrival times.
  \item Find the smallest unassigned time $t_{\min}$. This time must be the start of
        some route, so we try all routes beginning at $t_{\min}$.
  \item For each interval $d \in [1, 59]$, check the maximum number of evenly spaced
        arrivals $t_{\min}, t_{\min}+d, t_{\min}+2d, \ldots$ that are all available.
        For each valid count $c \ge 2$ (tried largest first to cover more arrivals per route),
        subtract the route from $\mathtt{cnt}$ and recurse.
  \item Prune: if the current number of routes already equals the best found, stop.
\end{enumerate}

\subsection{Correctness}

The key invariant is that the smallest unassigned time must start some route.
By trying all possible intervals and counts for that starting time, we enumerate
all valid decompositions. The pruning condition ensures we find a minimum-route solution.

\section{C++ Solution}

\begin{lstlisting}
#include <cstdio>
#include <cstring>
#include <vector>
#include <algorithm>
using namespace std;

int n;
int cnt[60]; // available count for each arrival time
int bestRoutes;

struct Route {
    int start, interval, count;
};
vector<Route> bestSolution, currentSolution;

bool canUse(int s, int d, int c) {
    for (int i = 0; i < c; i++) {
        int t = s + i * d;
        if (t > 59 || cnt[t] <= 0) return false;
    }
    return true;
}

void applyRoute(int s, int d, int c, int delta) {
    for (int i = 0; i < c; i++)
        cnt[s + i * d] += delta;
}

int remaining() {
    int r = 0;
    for (int i = 0; i < 60; i++) r += cnt[i];
    return r;
}

void solve() {
    if (remaining() == 0) {
        if ((int)currentSolution.size() < bestRoutes) {
            bestRoutes = (int)currentSolution.size();
            bestSolution = currentSolution;
        }
        return;
    }

    // Prune: cannot improve on current best
    if ((int)currentSolution.size() + 1 >= bestRoutes) return;

    // Find first unassigned time (must be the start of some route)
    int first = -1;
    for (int i = 0; i < 60; i++)
        if (cnt[i] > 0) { first = i; break; }

    // Try all routes starting at 'first'
    for (int d = 1; d <= 59; d++) {
        // Count maximum consecutive multiples present
        int maxc = 0;
        for (int t = first; t <= 59; t += d) {
            if (cnt[t] > 0) maxc++;
            else break;
        }
        if (maxc < 2) continue;

        // Try from largest count downward (greedy: cover more with fewer routes)
        for (int c = maxc; c >= 2; c--) {
            if (canUse(first, d, c)) {
                applyRoute(first, d, c, -1);
                currentSolution.push_back({first, d, c});
                solve();
                currentSolution.pop_back();
                applyRoute(first, d, c, +1);
            }
        }
    }
}

int main() {
    scanf("%d", &n);
    memset(cnt, 0, sizeof(cnt));
    for (int i = 0; i < n; i++) {
        int t;
        scanf("%d", &t);
        cnt[t]++;
    }

    bestRoutes = n; // upper bound
    solve();

    printf("%d\n", bestRoutes);
    for (auto& r : bestSolution)
        printf("Start: %d, Interval: %d, Count: %d\n",
               r.start, r.interval, r.count);

    return 0;
}
\end{lstlisting}

\section{Complexity Analysis}

\begin{itemize}
  \item \textbf{Time complexity:} Exponential in the worst case, but heavily pruned.
        At each recursion level, we fix the smallest unassigned time and try $O(59)$
        intervals, each with $O(59)$ possible counts. The pruning bound on the
        number of routes limits the recursion depth. In practice, for $n \le 300$
        and times $\le 59$, the search terminates quickly.
  \item \textbf{Space complexity:} $O(n)$ for the route stack and the frequency array.
\end{itemize}

\end{document}
